\appendixitem{李群的射影极限}

引理 1. 设 \((\mathrm{G}_{\alpha}, f_{\alpha \beta})\) 是相对于滤指标集 I 的拓扑群射影系统，G 是其极限。假设典则映射 \(f_{\alpha}: \mathrm{G} \to \mathrm{G}_{\alpha}\) 都是满射。

\begin{enumerate}
\def\labelenumi{\alph{enumi})}
\item
  子群 \(\overline{\mathrm{D}(\mathrm{G}_{\alpha})}\)（相应地，\(\mathrm{C}(\mathrm{G}_{\alpha})\)，相应地，\(\mathrm{C}(\mathrm{G}_{\alpha})_0\)）构成 \(\mathrm{G}_{\alpha}\) 的子集射影系统。
\item
  有 \(\overline{\mathrm{D}(\mathrm{G})}=\varprojlim_{\alpha}\overline{\mathrm{D}(\mathrm{G}_{\alpha})}\) 和 \(\mathrm{C}(\mathrm{G})=\varprojlim_{\alpha}\mathrm{C}(\mathrm{G}_{\alpha})\)。
\item
  若 \(\mathrm{G}_{\alpha}\) 对所有 \(\alpha \in \mathrm{I}\) 都是紧的，则 \(\mathrm{C}(\mathrm{G})_0 = \varprojlim_{\alpha} \mathrm{C}(\mathrm{G}_{\alpha})_0\)。
\end{enumerate}

设 \(\alpha, \beta\) 是 I 中满足 \(\alpha \leq \beta\) 的两个元素。则 \(f_{\alpha \beta}(\mathrm{D}(\mathrm{G}_{\beta})) \subset \mathrm{D}(\mathrm{G}_{\alpha})\)，并且 \(f_{\alpha \beta}(\mathrm{C}(\mathrm{G}_{\beta})) \subset \mathrm{C}(\mathrm{G}_{\alpha})\)，因为 \(f_{\alpha \beta}\) 是满射；由于 \(f_{\alpha \beta}\) 连续，遂有 \(f_{\alpha \beta}(\overline{\mathrm{D}(\mathrm{G}_{\beta})}) \subset \overline{\mathrm{D}(\mathrm{G}_{\alpha})}\) 以及 \(f_{\alpha \beta}(\mathrm{C}(\mathrm{G}_{\beta})_0) \subset \mathrm{C}(\mathrm{G}_{\alpha})_0\)，故得 \(a)\)。由于 \(f_{\alpha}\) 是满射，\(f_{\alpha}(\mathrm{D}(\mathrm{G})) = \mathrm{D}(\mathrm{G}_{\alpha})\)（Algebra, Chap. I, §6, no. 2, Prop. 6），所以 \(\overline{\mathrm{D}(\mathrm{G})} = \varprojlim \overline{\mathrm{D}(\mathrm{G}_{\alpha})}\)（General Topology, Chap. I, §4, no. 4, Cor. of Prop. 9）。\(f_{\alpha}\) 的满射性还推出包含关系 \(f_{\alpha}(\mathrm{C}(\mathrm{G})) \subset \mathrm{C}(\mathrm{G}_{\alpha})\)，从而 \(\mathrm{C}(\mathrm{G}) \subset \varprojlim \mathrm{C}(\mathrm{G}_{\alpha})\)；反向包含显然。最后，断言 \(c)\) 由 \(b)\) 和 General Topology, Chap. III, §7, no. 2, Prop. 4 推出。

引理 2. 设 \((\mathrm{S}_a)_{a\in \mathrm{A}},(\mathrm{T}_b)_{b\in \mathrm{B}}\) 是两个有限的概单单连通李群族（Chap. III, § 9, no. 8, Def. 3），\(u:\prod_{a\in \mathrm{A}}\mathrm{S}_a\to \prod_{b\in \mathrm{B}}\mathrm{T}_b\) 是一个满射态射。则存在单射 \(l: \mathrm{B} \to \mathrm{A}\) 以及同构 \(u_b: \mathrm{S}_{l(b)} \to \mathrm{T}_b\)（\(b \in \mathrm{B}\)），使得 \(u((s_a)_{a \in \mathrm{A}}) = (u_b(s_{l(b)}))_{b \in \mathrm{B}}\)，对每个元素 \((s_a)_{a \in \mathrm{A}}\)（属于 \(\prod_{a \in \mathrm{A}} \mathrm{S}_a\)）都成立。

记 \(\mathfrak{s}_a\)（相应地，\(\mathfrak{t}_b\)）为 \(S_a\)（相应地，\(T_b\)）的李代数，其中 \(a \in A\)（相应地，\(b \in B\)），并考虑同态 \(L(u): \prod_{a \in A} \mathfrak{s}_a \to \prod_{b \in B} \mathfrak{t}_b\)。其核是半单李代数 \(\prod_{a \in A} \mathfrak{s}_a\) 的一个理想，因而形如 \(\prod_{a \in A''} \mathfrak{s}_a\)，其中 \(A'' \subset A\)（Chap. I, §6, no. 2, Cor. 1）。置 \(A' = A - A''\)。限制 \(L(u)\) 后得到同构 \(f: \prod_{a \in A'} \mathfrak{s}_a \to \prod_{b \in B} \mathfrak{t}_b\)。由 loc. cit.，对于所有 \(a \in A'\)，理想 \(f(\mathfrak{s}_a)\) 等于某个 \(\mathfrak{t}_b\)；因此存在双射 \(l: B \to A'\)，使得 \(f(\mathfrak{s}_{l(b)}) = \mathfrak{t}_b\) 对于 \(b \in B\) 成立，并且 \(f\) 诱导出同构 \(f_b: \mathfrak{s}_{l(b)} \to \mathfrak{t}_b\)。由于群 \(S_a\) 和 \(T_b\) 单连通，存在同构 \(u_b: S_{l(b)} \to T_b\)，使得 \(L(u_b) = f_b\) 对于 \(b \in B\) 成立（Chap. III, §6, no. 3, Th. 3）。

记 \(\tilde{u}:\prod_{a\in \mathrm{A}}\mathrm{S}_a\to \prod_{b\in \mathrm{B}}\mathrm{T}_b\) 为由 \(\tilde{u} ((s_a)_{a\in \mathrm{A}}) = (u_b(s_{l(b)}))_{b\in \mathrm{B}}\) 定义的态射。按构造，\(\mathrm{L}(\tilde{u}) = f = \mathrm{L}(u)\)，所以 \(\tilde{u} = u\)，引理得证。

引理 3. 在引理 1 的假设下，假设 \(G_{\alpha}\) 是单连通紧李群。则拓扑群 G 同构于一族概单单连通紧李群的乘积。

对于所有 \(\alpha \in \mathrm{I}\)，群 \(\mathrm{G}_{\alpha}\) 是有限族概单单连通子群 \((\mathrm{S}_{\alpha}^{\lambda})_{\lambda \in \mathrm{L}_{\alpha}}\) 的直积（Chap. III, §9, no. 8, Prop. 28）。设 \(\beta \in \mathrm{I}\)，\(\beta \geq \alpha\)。由引理 2，存在映射 \(l_{\beta \alpha}: \mathrm{L}_{\alpha} \to \mathrm{L}_{\beta}\)，使得 \(f_{\alpha \beta}(\mathrm{S}_{\beta}^{l_{\beta \alpha}(\lambda)}) = \mathrm{S}_{\alpha}^{\lambda}\) 对于 \(\lambda \in \mathrm{L}_{\alpha}\) 成立。有 \(l_{\gamma \beta} \circ l_{\beta \alpha} = l_{\gamma \alpha}\)，其中 \(\alpha \leq \beta \leq \gamma\)，所以 \((\mathrm{L}_{\alpha}, l_{\beta \alpha})\) 是相对于 I 的集合归纳系统。令 L 为其极限；由于映射 \(l_{\beta \alpha}\) 是单射，可以将 \(\mathrm{L}_{\alpha}\) 识别为 L 的子集，从而 \(\mathrm{L} = \bigcup_{\alpha \in \mathrm{I}} \mathrm{L}_{\alpha}\)。

设 \(\lambda \in L\)。置 \(S_{\alpha}^{\lambda} = \{1\}\)，当 \(\lambda \notin L_{\alpha}\) 时，并记 \(\varphi_{\alpha \beta}^{\lambda}: S_{\beta}^{\lambda} \to S_{\alpha}^{\lambda}\) 为由 \(f_{\alpha \beta}\) 诱导的态射；这给出拓扑群射影系统 \((S_{\alpha}^{\lambda}, \varphi_{\alpha \beta}^{\lambda})\)，其极限同构于 \(S_{\lambda}\)。拓扑群的典则同态

\[
\lim_{\substack{\longleftarrow \\ \alpha \in I}}\left(\prod_{\lambda \in L}S_{\alpha}^{\lambda}\right)\to \prod_{\lambda \in L}\left(\lim_{\substack{\longleftarrow \\ \alpha \in I}}S_{\alpha}^{\lambda}\right)
\]

是双射（Theory of Sets, Chap. III, §7, no. 3, Cor. 2）；由于有关群都是紧的，它因而是同构。前一个群可与 G 识别，后一个可与 \(S_{\lambda}\) 的乘积识别，故引理得证。

